% f_(a,b)(x)=a tanh(bx)，x*=1、y=.9、半平方损失、eta=.3，t=0,...,6。
% 三面板绘图区均为44mm x 34mm；坐标变量、单位和范围各自标明。
% 真实与一阶输出在同一条真实参数轨迹上求值，不是任意参数切线投影。
\begin{tikzpicture}[foundation picture]
 \path[use as bounding box] (0,-11) rectangle (169,63);
 \foreach \origin in {10,65,121}{
  \draw[foundation axis] (\origin,13)--(\origin+44,13);
  \draw[foundation axis] (\origin,13)--(\origin,47);
 }
 \node[above] at (10,47) {$b$};
 \node at (54,3) {$a$};
 \foreach \v in {1,1.05,1.1}{\node[below] at ({10+44*(\v-1)/.13},12) {$\v$};}
 \foreach \v in {1,1.03,1.06}{\node[left] at (8.5,{13+34*(\v-1)/.08}) {$\v$};}
 \node[above] at (65,47) {$f(x_*)$};
 \foreach \v/\vlabel in {.75/0.75,.80/0.80,.85/0.85,.90/0.90}{
  \node[left] at (63.5,{13+34*(\v-.75)/.15}) {$\vlabel$};
 }
 \foreach \origin in {65,121}{
  \foreach \t in {0,2,4,6}{\node[below] at ({\origin+44*\t/6},12) {$\t$};}
  \node at (\origin+22,3) {更新次数$t$};
 }
 \node at (143,53) {$|f-f^{\mathrm{lin}}|$};
 \foreach \err/\label in {0/0,.001/0.001,.002/0.002}{
  \node[left] at (119.5,{13+34*\err/.002}) {$\label$};
 }
 % 以下七行由已声明的梯度下降递推计算，数据不因版式修改而改变。
 \foreach \t/\a/\b/\real/\linear/\err in {
  0/1.000000000000/1.000000000000/0.761594155956/0.761594155956/0.000000000000,
  1/1.031622724592/1.017438070968/0.793133202559/0.793001380575/0.000131821984,
  2/1.056271155636/1.030962446940/0.817864478461/0.817453372505/0.000411105955,
  3/1.075350269497/1.041385532109/0.837089867821/0.836361342454/0.000728525366,
  4/1.090041696762/1.049382615479/0.851939763898/0.850908817425/0.001030946472,
  5/1.101310373819/1.055498679800/0.863355699177/0.862059566104/0.001296133073,
  6/1.109928398322/1.060165283821/0.872100833020/0.870582857153/0.001517975868
 }{
  \coordinate (parameter\t) at ({10+44*(\a-1)/.13},{13+34*(\b-1)/.08});
  \coordinate (real\t) at ({65+44*\t/6},{13+34*(\real-.75)/.15});
  \coordinate (linear\t) at ({65+44*\t/6},{13+34*(\linear-.75)/.15});
  \coordinate (error\t) at ({121+44*\t/6},{13+34*\err/.002});
 }
 \draw[FigureYellow,line width=1.1pt] plot coordinates {
  (parameter0) (parameter1) (parameter2) (parameter3) (parameter4) (parameter5) (parameter6)};
 \draw[FigureBlue,line width=1.1pt] plot coordinates {
  (real0) (real1) (real2) (real3) (real4) (real5) (real6)};
 \draw[FigureRed,line width=1.1pt,dashed] plot coordinates {
  (linear0) (linear1) (linear2) (linear3) (linear4) (linear5) (linear6)};
 \draw[FigureYellow,line width=1.1pt] plot coordinates {
  (error0) (error1) (error2) (error3) (error4) (error5) (error6)};
 \foreach \t in {0,...,6}{
  \fill[FigureYellow] (parameter\t) circle (.85);
  \fill[FigureBlue] (real\t) circle (.75);
  \draw[FigureRed,fill=none,line width=.8pt] (linear\t) +(-.9,-.9) rectangle +(.9,.9);
  \fill[FigureYellow] (error\t) circle (.75);
 }
 \draw[FigureYellow,fill=white,line width=1pt] (parameter0) circle (1.2);
 \node at (19,3) {$\bm\theta^{(0)}$};
 \node at (44,44) {$\bm\theta^{(6)}$};
 \draw[FigureBlue,line width=1.1pt] (65,59)--(72,59);
 \fill[FigureBlue] (68.5,59) circle (.75);
 \node[anchor=west] at (74,59) {真实};
 \draw[FigureRed,line width=1.1pt,dashed] (94,59)--(101,59);
 \draw[FigureRed,fill=none,line width=.8pt] (97.5,59) +(-.9,-.9) rectangle +(.9,.9);
 \node[anchor=west] at (103,59) {一阶};
 \node at (84.5,-7) {$f_{\bm\theta}^{\mathrm{lin}}(x)=f_0(x)+\langle\bm\varphi_0(x),\bm\theta-\bm\theta^{(0)}\rangle$};
\end{tikzpicture}%
